\section{Introduction}\label{sec:intro}

Zero-knowledge (ZK) proofs now guard real money on public blockchains.
Privacy pools release deposits to whoever proves membership, rollups settle
state on a proof of correct execution, identity systems mint credentials to
the holder of a proved attribute, and zkVM markets pay provers for proved
computations. In each case a smart contract, the \emph{consumer}, accepts a
proof and then acts: it transfers, mints, records or grants. A great deal of
work makes sure that proofs are sound, that circuits are fully constrained,
and that verifiers are generated correctly~\cite{sok-snark-bugs,picus,circomspect,zkap}.
Much less attention goes to a simpler question: does the contract's action
follow from what was proved?

It often does not. A proof certifies a \emph{statement}, the public inputs
and what the circuit fixes about them. The consumer performs an \emph{effect},
which depends on a recipient, an amount, a fee, a state root, a replay key,
and a verification key. When an effect value is not fixed by the statement,
whoever submits the proof chooses it. Valid proofs are public once broadcast,
so that submitter can be an attacker who copies a victim's proof and changes
one argument. Figure~\ref{fig:motivating} shows the shape of such a flaw. We
call it a \emph{binding flaw}: the proof is valid, the verifier is correct,
and the effect is still unauthorized.

Binding flaws are not hypothetical. In June 2026 an attacker drained \$2.28M
from the deprecated Aztec Connect rollup contract without breaking any
cryptography. Each of its batches carried a valid proof, but the contract
settled only as many rows as a calldata field, \texttt{numRealTxs}, declared,
and the proof did not bind that field~\cite{aztec-connect-hack} (P1 in our
taxonomy). ZKSwap's bridge lost \$5M in 2025 to an exit verifier that
returned \emph{true} before checking anything~\cite{zkswap-hack}, a
verifier-integration flaw (P4). These two core-class incidents account for
about \$7.3M of actual losses. A third incident lies outside our scope: in
2026 two privacy pools put about \$1.4--2.3M at risk, most of it later
rescued, because their verification keys skipped the second phase of the
trusted setup~\cite{groth16-setup-exploit,rekt-foom}, a setup failure rather
than a binding flaw. In a public record of ZK incidents from 2023 to 2026,
every circuit-soundness bug in an on-chain application was caught before
exploitation, while each quantified loss we could trace sits at the verifier
or at the contract that consumes it~\cite{zkp-hacks-record}.
% src: rekt.news/aztec-connect-rekt ($2.28M; 14 valid processRollup calls; numRealTxs not bound, BlockSec); blockaid ZKSwap ($5M, 2025-07-09, verifyExitProof returns true); zksecurity 2026-02-27 (Foom ~$1.4-1.5M, Veil ~$5K) + rekt the-unfinished-proof ($2.26M gross); floor 2.28+5+1.4 = 8.68

\begin{figure}[t]
\begin{lstlisting}[language=Java]
// Circuit: public [root, nullifierHash]; proves membership
function withdraw(bytes proof, uint root, uint nfHash,
                  address recipient, uint fee) external {
  require(isKnownRoot(root) && !spent[nfHash]);
  require(verifier.verify(proof, [root, nfHash]));
  spent[nfHash] = true;
  token.transfer(recipient, DENOM - fee);   // not in the statement
  token.transfer(msg.sender, fee);          // not in the statement
}
\end{lstlisting}
\caption{\DEMO{A binding flaw of the kind auditors report (simplified).}
The proof fixes the root and nullifier but not the recipient or the fee. Any
party that sees the transaction in the mempool can resubmit the proof with its
own recipient. \TODO{F1}{replace with a real, already-fixed audited case (dev split), with circuit, consumer and attack in three panels}}
\label{fig:motivating}
\end{figure}

\para{Binding flaws are common and severe, and today's tools are blind
to them}
We collected 8,578 public audit reports and labeled 264 candidate
findings about proof use. 87 of them are core binding flaws (103 with
adjacent proof-system bugs), spread over 43 of the 68 project families, and
78\% of those with a severity rating are Critical or High; 55 of the 311
in-scope reports contain at least one (Section~\ref{sec:study}). Yet, on the
audited flaws of our test set, the Solidity analyzers Slither and Mythril
report none at the right location and mechanism, on the 8 of 14 targets they
could analyze. The circuit analyzers circomspect and Picus cannot analyze most
targets, because most binding flaws are not in a circuit. A property-based
fuzzer with generic properties finds none of them
(Section~\ref{sec:eval}). The reason is structural. A binding flaw breaks no
rule of the contract and no constraint of the circuit. It breaks an
\emph{intended correspondence} between the two, which neither artifact states.

\para{This is not input validation}
Input validation checks a value against the program's own rules; binding
checks it against a statement proven elsewhere, by a different component. A
recipient that passes every contract-level check is still unbound if the
proof did not commit to it, which is why generic analyses miss it and why a
fuzzer finds these flaws once given the binding specification: 5 of 13,
against 0 without it.

The problem has been observed before. Chaliasos et al.\ report 13
integration-layer bugs among SNARK vulnerabilities, 9 of them unchecked public
inputs, and note that no existing defense covers that layer~\cite{sok-snark-bugs};
bug trackers document single instances such as nullifiers aliased modulo the
field order. We make the property precise, measure it at a larger scale
\DEMO{(49 of our findings fall in their integration layer, 3.8 times their
count, and one of our two most common patterns has no category of its own there)}, and
build the detector that this line of work calls for. The two-ledger view
names what a taxonomy of bug locations does not: the missing equality between
a statement row and an effect row, each correct on its own
(Section~\ref{sec:eval-misses}).

\para{This work}
We first make the property precise. We describe a consumer by two ledgers: a
\emph{statement ledger} of what the accepted proof fixes, and an \emph{effect
ledger} of every value that decides who, how much, what, which state and how
often. A binding flaw is an effect value that the attacker controls and that
the statement does not determine. Its forms mirror well-known signature
misuse: unsigned fields, missing domain separation, nonce reuse, weak or
unpinned keys, and format confusion.

We then build \tool{}, which makes an LLM construct both ledgers explicitly,
row by row, under a fixed procedure derived from development data only,
query deterministic tools for facts, and \DEMO{write and run attack tests on a
local chain, which it attaches to each warning.}
On 22 test-set audited flaws a first version of \tool{} finds 6 at the
exact location and mechanism (9 at the mechanism level), all of them in
Solidity consumers, and 22 of 27 injected mutants in three small hosts. The
existing analyzers find none of the flaws they could analyze. Tracing every miss shows that the brain read the flawed
code each time and then judged it wrongly, often about what a zkVM guest or
circuit commits. The final version therefore adds a knowledge base per proof
system and a second look that challenges every lenient verdict.
\DEMO{We evaluate the two additions together. Frozen before any held-out run,
the final version finds 10 of 14 held-out audited flaws (71\%, median of
three runs), against 2 of 7 on the same check set for the first version and
4 of 14 for the same LLM without the procedure under equal runs. On 37 new
injected flaws in five larger hosts it finds 31, including 8 of the 11
circuit- or guest-side ones. On 25 recently audited applications it found
three previously unknown binding flaws, each confirmed with a real proof on a
local chain and reported to the maintainers.}

\para{Contributions}
\begin{itemize}
\item \textbf{A definition and taxonomy} of proof-to-effect binding flaws:
two ledgers, three conditions, and five patterns with signature-misuse
analogues (Section~\ref{sec:definition}).
\item \textbf{To our knowledge, the first measurement study} of the class: 8,578 audit reports,
264 labeled findings, and 87 core binding flaws, with prevalence, severity,
patterns, and trends (Section~\ref{sec:study}).
\item \textbf{\tool{}}, a detector that combines an LLM with a ledger-based
skill procedure, deterministic tools, per-proof-system knowledge, a
second-look stage, and an exploration loop that attaches executed attack tests
to its warnings (Section~\ref{sec:design}).
\item \textbf{An evaluation} against static, circuit, fuzzing and LLM
baselines on real audited flaws and on a new binding-mutant benchmark,
\DEMO{including three new flaws confirmed with real-proof tests} (Section~\ref{sec:eval}).
We release the dataset, the benchmark and the tool.
\end{itemize}
